\chapter{Conclusão}

O MiniLang / MiniLexer atingiu uma implementação funcional coerente com o escopo definido para um analisador léxico. O projeto contém seis Expressões Regulares principais, seis construções de AFN$\varepsilon$, simulador genérico, lexer com regras de prioridade e maior lexema, tratamento de erros e suíte de testes integrada.

A verificação final registrou 135 testes automatizados aprovados. As seis validações cruzadas AFN$\varepsilon$ $\times$ regex Python não apresentaram divergências nos alfabetos reduzidos definidos, e a validação estrutural confirmou 6/6 implementações Python e 6/6 arquivos JFLAP contra as tabelas canônicas.

A principal contribuição acadêmica da implementação está na manutenção da correspondência entre a descrição formal e o comportamento computacional. Cada ER é apresentada com sua finalidade, alfabeto, linguagem, notação formal, sintaxe Python, AFN$\varepsilon$, casos dirigidos e limitações, enquanto os testes de integração demonstram como essas regras funcionam em conjunto no MiniLexer.

Os artefatos finais --- código-fonte, testes, diagramas dos AFN$\varepsilon$, relatório, apresentação e registro de contribuições --- estão disponíveis no repositório do projeto.
